% Deutsche Parallelfassung zu why/vision.tex

\chapter{Die AIM-These}

\section{Von der Folge zur Ingenieurumgebung}

Die vorangegangenen Kapitel haben keinen Mangel an Repräsentationen
offengelegt. Sie haben einen Mangel an Kontinuität zwischen ihnen sichtbar
gemacht. Die Krümmungsänderung kann gezeichnet, verortet, ausgetauscht,
gemessen, geprüft und erörtert werden; jede Operation sieht jedoch einen
anderen Teil ihrer Bedeutung. AIM antwortet, indem es eine Ingenieurumgebung
um das Alignment und die Abhängigkeiten organisiert, durch die eine Handlung
Folgen erhält.

In dieser Umgebung kann der Ingenieur eine zusammenhängende Folge von Fragen
stellen: Welches Alignment wird geändert? Welches konstruktive Gesetz ändert
sich? Welche Zustände folgen daraus? Welche Qualifikatoren und Beobachtungen
sind relevant? Welches Wissen und welche Nebenbedingungen sind anwendbar?
Welche Evidenz stützt Annahme oder Ablehnung? Das System beantwortet nicht alle
diese Fragen auf dieselbe Weise. Es bewahrt ihre Beziehungen, damit jedes
verantwortliche Werkzeug oder jede verantwortliche Person die jeweils
angemessene Frage beantworten kann.

\begin{proposition}[AIM-These]
Alignment-Based Information Modelling ist eine Alignment-spezifische
Ingenieurumgebung, in der dauerhafte Objektidentität, konstruktive Definition,
Realisierung, Repräsentation, Beobachtung, anwendbares Wissen, Bewertung und
Entscheidung im Handlungsmoment getrennt, aber ausdrücklich aufeinander bezogen
bleiben.
\end{proposition}

\section{Von der Krümmung zur Folge}

Die fortlaufende Änderung veranschaulicht die erforderliche Kette. Der
geordnete Krümmungsverlauf gehört zur intrinsischen Konstruktion des Alignments.
Unter einer metrischen Interpretation erzeugt er die Lagegeometrie; mit
weiterem Zustand und weiteren Qualifikatoren unterstützt er räumliche,
dynamische und regelbasierte Bewertung. Unterschiedliche Repräsentationen
legen unterschiedliche Teile dieser Kette offen. Keine besitzt die gesamte
Kette.

\begin{figure}[H]
\centering
% Deutsche Parallelfassung zu curvature_consequence_chain.tex
\begin{tikzpicture}[
  x=1cm,
  y=1cm,
  >=Stealth,
  flow/.style={->, semithick, draw=black!65},
  stage/.style={draw=black!55, rounded corners=2pt, fill=black!2,
    minimum width=3.05cm, minimum height=2.55cm, inner sep=5pt},
  label/.style={font=\scriptsize\bfseries, align=center},
  note/.style={font=\scriptsize, align=center, text=black!75}
]
  \node[stage] (kappa) at (0,0) {};
  \node[label, anchor=north] at ([yshift=-3pt]kappa.north)
    {Krümmungsband\\$\kappa(s)$};
  \draw[->, thin] ([xshift=-1.15cm,yshift=-0.65cm]kappa.center) -- ++(2.35,0)
    node[right, font=\scriptsize] {$s$};
  \draw[->, thin] ([xshift=-1.05cm,yshift=-0.85cm]kappa.center) -- ++(0,1.25);
  \draw[very thick, black!75]
    plot[smooth] coordinates {
      ([xshift=-1.00cm,yshift=-0.48cm]kappa.center)
      ([xshift=-0.55cm,yshift=-0.42cm]kappa.center)
      ([xshift=-0.10cm,yshift=-0.05cm]kappa.center)
      ([xshift=0.38cm,yshift=0.35cm]kappa.center)
      ([xshift=1.00cm,yshift=0.39cm]kappa.center)};
  \node[note, anchor=south] at ([yshift=3pt]kappa.south)
    {intrinsische Konstruktion};

  \node[stage, right=0.55cm of kappa] (plan) {};
  \node[label, anchor=north] at ([yshift=-3pt]plan.north)
    {Lagegeometrie\\$\gamma(s)$};
  \draw[very thick, black!75]
    ([xshift=-1.10cm,yshift=-0.40cm]plan.center)
    .. controls ([xshift=-0.25cm,yshift=-0.40cm]plan.center)
      and ([xshift=0.15cm,yshift=0.48cm]plan.center)
    .. ([xshift=1.10cm,yshift=0.48cm]plan.center);
  \draw[dashed, black!45]
    ([xshift=-1.10cm,yshift=-0.18cm]plan.center)
    .. controls ([xshift=-0.28cm,yshift=-0.18cm]plan.center)
      and ([xshift=0.18cm,yshift=0.70cm]plan.center)
    .. ([xshift=1.10cm,yshift=0.70cm]plan.center);
  \node[note, anchor=south] at ([yshift=3pt]plan.south)
    {metrische Interpretation};

  \node[stage, right=0.55cm of plan] (context) {};
  \node[label, anchor=north] at ([yshift=-3pt]context.north)
    {Anwendbarkeits-\\qualifikatoren};
  \draw[black!25] ([xshift=-1.10cm,yshift=-0.45cm]context.center)
    grid[xstep=.45cm,ystep=.38cm]
    ([xshift=1.10cm,yshift=.70cm]context.center);
  \draw[very thick, black!75]
    ([xshift=-1.05cm,yshift=-0.35cm]context.center)
    .. controls ([xshift=-0.25cm,yshift=-0.30cm]context.center)
      and ([xshift=0.15cm,yshift=0.48cm]context.center)
    .. ([xshift=1.05cm,yshift=0.52cm]context.center);
  \fill[black!55] ([xshift=.48cm,yshift=.15cm]context.center) circle (1.6pt);
  \node[note, anchor=south] at ([yshift=3pt]context.south)
    {Realität, Regeln, Beobachtungen};

  \node[stage, right=0.55cm of context] (use) {};
  \node[label, anchor=north] at ([yshift=-3pt]use.north)
    {Zweckgebundene\\Darstellungen};
  \node[draw=black!45, rounded corners=1pt, font=\scriptsize,
    minimum width=.72cm] at ([xshift=-.78cm,yshift=.25cm]use.center) {GIS};
  \node[draw=black!45, rounded corners=1pt, font=\scriptsize,
    minimum width=.72cm] at ([yshift=.25cm]use.center) {CAD};
  \node[draw=black!45, rounded corners=1pt, font=\scriptsize,
    minimum width=.72cm] at ([xshift=.78cm,yshift=.25cm]use.center) {BIM};
  \node[note] at ([yshift=-.32cm]use.center)
    {Bewertung\;\textbullet\; Folge};
  \node[note, anchor=south] at ([yshift=3pt]use.south)
    {Evidenz für menschliche Entscheidung};

  \draw[flow] (kappa) -- (plan);
  \draw[flow] (plan) -- (context);
  \draw[flow] (context) -- (use);
\end{tikzpicture}

\caption{Eine Krümmungsänderung, mehrere notwendige Interpretationen. Das
Krümmungsband erzeugt Lagegeometrie; getrennte Anwendbarkeitsqualifikatoren
beziehen diese Geometrie auf Realität und Regeln; zweckgebundene
Repräsentationen legen ausgewählte Informationen offen; Bewertungen zeigen
Folgen. Die Pfeile drücken Abhängigkeit aus, nicht Ownership oder automatische
Entscheidungsautorität.}
\label{fig:curvature-consequence-chain}
\end{figure}

Abbildung~\ref{fig:curvature-consequence-chain} ist bewusst keine
Softwarepipeline, sondern eine Argumentationsabbildung. Von links nach rechts
gelesen verfolgt sie die Folge einer Änderung; von rechts nach links gelesen
fragt sie, welche Provenienz und Konstruktion zur Begründung eines angezeigten
Ergebnisses erforderlich sind. Eine GIS-Sicht kann räumliche Qualifikationen
liefern, eine CAD-Sicht genaue Geometrie offenlegen, eine BIM-Sicht
Austauschobjekte tragen und eine Beobachtung über das Gleis berichten.
Vergleichbar werden ihre Aussagen erst, wenn ihre Beziehung zu demselben
Engineering Object und zum maßgeblichen Zustand bekannt ist.

Für den ebenen Bogenlängenparameter \(s\) ist die erste konstruktive
Abhängigkeit kompakt:
\[
  \theta'(s)=\kappa(s),
  \qquad
  \gamma'(s)=\bigl(\cos\theta(s),\sin\theta(s)\bigr).
\]
Das Krümmungsgesetz bestimmt die Entwicklung des Richtungswinkels, und der
Richtungswinkel bestimmt den ebenen Pfad. Literatur zu
Eisenbahn-Übergangsbögen und Entwurfsnormen verbinden denselben
Krümmungsverlauf zusammen mit Überhöhung und Geschwindigkeit mit
Beschleunigung, Komfort und Zulässigkeit, nicht allein mit sichtbarer Form
\cite{Kufver1997MathematicalDescription,Kufver2000RailwayAlignment,EN13803_2017,BrustadDalmo2020Review}.

\section{Wissen im Handlungsmoment}

Wissen anwendbar zu machen bedeutet nicht, jede Regel in jedes Modell zu
kopieren. Es bedeutet, dass eine Änderung auf Wissen bezogen werden kann,
dessen Geltungsbereich, Provenienz, Anwendbarkeit und Autorität ausdrücklich
sind. Wenn sich das Krümmungsband ändert, kann die Umgebung betroffene Zustände
bestimmen und die relevanten Bewertungen aufrufen. Sie kann zeigen, welche
Beobachtung einen Vergleich stützt und welche Nebenbedingung verletzt ist. Sie
kann bewahren, warum das Ergebnis hier anwendbar ist.

Dies verlangt mehr, als Attribute an eine Kurve anzuhängen. Attribute können
Werte speichern; Anwendbarkeit erfordert Beziehungen zwischen Objekt, Zustand,
Qualifikatoren, Wissen und Operation. Wegen dieser Beziehungen sind formale
Grundlagen erforderlich.

\section{Verantwortlichkeiten, die nicht zusammenfallen}

AIM macht weder ein Modell noch ein Programm zu einer universellen Autorität.

\begin{description}[style=nextline]
\item[Repräsentationen]
GIS, CAD, BIM, Tabellen, Diagramme und Austauschdateien wählen Informationen
für einen Zweck aus. Sie kommunizieren und unterstützen Handlungen, erzeugen
aber nicht allein dadurch ingenieurtechnische Identität oder Autorität, dass
sie einen Wert speichern.

\item[Modelle und Operatoren]
Konstruktive und Zustandsmodelle machen Abhängigkeiten ausdrücklich.
Operatoren leiten Zustände unter benannten Annahmen ab, transformieren,
vergleichen und bewerten sie.

\item[Berechnung]
Numerische Verfahren können Zustände rekonstruieren, Folgen fortpflanzen,
Nebenbedingungen prüfen und Lösungskandidaten erzeugen. Das Ergebnis eines
Solvers ist Evidenz, nicht Annahme.

\item[Ingenieure und verantwortliche Autoritäten]
Menschen bestimmen Anwendbarkeit, interpretieren Evidenz, verantworten
Ingenieurwissen und entscheiden innerhalb ihrer Autorität, ob eine
vorgeschlagene Änderung angenommen wird.
\end{description}

Die Referenzanwendung ist ein Realisierungshorizont für diese Trennung. Sie
kann ein Krümmungsband bearbeitbar machen, eine geometrische Sicht ableiten,
Qualifikationen zuordnen und Bewertungsergebnisse in einer Interaktion
sichtbar machen. Ihr Wert für die Thesis liegt nicht in einer
Produktdemonstration, sondern in einem Test: Können die Begriffe getrennt
bleiben, wenn ein Ingenieur tatsächlich etwas ändert?

\section{Warum Alignment-spezifisch?}

Allgemeinheit ist erst nützlich, wenn die bestimmende Abhängigkeit sichtbar
ist. Eisenbahnlage, Richtung, Überhöhung, fahrzeugbezogene Größen,
Stationierung, benachbarte Objekte und viele Ingenieurregeln werden entlang
oder relativ zu einem Alignment interpretiert. Krümmungsverlauf und geordnete
longitudinale Konstruktion bilden daher ein domänenspezifisches Rückgrat, das
ein allgemeiner Asset-Container oder ein raumbezogenes Feature allein nicht
bereitstellt.

Alignment-spezifisch bedeutet nicht Alignment-exklusiv. Ingenieurbauwerke,
Signale, Gelände, Flurstücke, Fahrzeuge und Projektorganisationen bleiben
Engineering Objects mit eigener Identität und eigenem Wissen. AIM beschreibt,
wie sie auf ein Alignment bezogen werden können, ohne in ihm aufzugehen.

\section{Die Schwelle zu den Grundlagen}

Die These erzeugt nun eine formale Verpflichtung. Wenn Identität einen Wechsel
der Repräsentation überstehen muss, benötigen Identität und Repräsentation
unterschiedliche formale Rollen. Wenn eine Beobachtung Evidenz statt Realität
ist, benötigen Beobachtung und physischer Zustand unterschiedliche
Beziehungen. Wenn Krümmung Geometrie erzeugt, benötigen das konstruktive Gesetz
und seine metrische Realisierung einen ausdrücklichen Operator. Wenn Bewertung
informiert, aber nicht entscheidet, dürfen Kandidat, Ergebnis und Entscheidung
nicht dasselbe undifferenzierte Statusfeld teilen.

Der nächste Teil entwickelt die minimale begriffliche und mathematische
Sprache, die erforderlich ist, um diese Unterscheidungen zu bewahren und
ausgewählte Abhängigkeiten berechenbar zu machen. Er beginnt die Argumentation
nicht neu. Er gibt ihr eine Form, die geprüft werden kann.

\begin{summary}
AIM verbindet eine Alignment-Änderung mit ihren Folgen, indem es ausdrückliche
Beziehungen zwischen Konstruktion, Zustand, Qualifikatoren, Evidenz, Wissen,
Bewertung und Autorität bewahrt. Die folgende formale Arbeit ist notwendig,
weil diese Beziehungen nicht zuverlässig allein aus Geometrie oder
Datencontainern zurückgewonnen werden können.
\end{summary}
